In this section we will discuss how we implemented a Mixed Integer Linear Program model able to solve the given Tournament Scheduling problem.

The main obstacle in the implementation of linear programs is the fact that the objective function and every constraint have to be linear in our variables; therefore we transformed logic propositions and non-linear functions with linear inequalities, introducing, if needed, some auxiliary variable as for the implementation of the objective function.

\subsection{Decision variables}
Before defining the variables, we compute the list $L$ of the possible match-ups, without caring about the order of the $(i,j)$ team tuples (we impose $i<j$ for unicity).
If $n=T$ is the number of teams, we have $|L|=\binom{n}{2} = \frac{n!}{2!(n-2)!}$, i.e. $\binom{n}{2}$ games to be played.


Our decision variables are then two 3-dimensional tensors of \textbf{boolean} values $Y_{w,p,(i,j)}$ and $H_{w,p,(i,j)}$, where
\begin{itemize}
    \item $Y_{w,p,(i,j)} = 1 \iff$ teams $i$ and $j$ are playing one against the other in week $w$ in period $p$;
    \item $H_{w,p,(i,j)} = 1 \iff$ team $i$ (the team with lowest index since $\forall (i,j) \in L$, $i<j$) is playing home in week $w$ in period $p$.
\end{itemize}

Clearly these variables have to be linked one with the other, and that's why we will introduce some variable-linking constraint to our model.


\subsection{Objective function}

The objective function is used only in the optimization version of the MIP
model. Its purpose is to balance the number of home and away games played by
each team.

For each team $t$, let $h_t$ be the number of home games played by team $t$.
Since each team plays exactly one match per week and there are $W = n - 1$
weeks, the number of away games is:

\[
a_t = W - h_t.
\]

The home/away imbalance of team $t$ is therefore:

\[
|h_t - a_t|.
\]

Since absolute values cannot be used directly in a linear objective function, we
introduce an auxiliary integer variable $b_t$ for each team $t$, representing
the imbalance of that team. We impose the following linear constraints:

\[
b_t \geq h_t - a_t \qquad \forall t,
\]

\[
b_t \geq a_t - h_t \qquad \forall t.
\]

Since the objective minimizes the sum of the imbalance variables, each $b_t$
will take the smallest value satisfying both inequalities, and therefore:

\[
b_t = |h_t - a_t|.
\]

The objective function of the optimization version is then:

\[
\min \sum_{t=1}^{n} b_t.
\]

In the decision version, no real objective function is optimized. A dummy
constant objective is used only to run the feasibility model through the MIP
solver.



% ======= CONSTRAINTS MIP =======
\subsection{Constraints}

\paragraph{Variable linking constraint}
As already mentioned, we have to link the variables and this has to be done by adding linear inequalities to our linear program.

The only constraint we have to add is the one that ensures that $H_{w,p,(i,j)}\neq1$ if teams $i$ and $j$ do not play against in week $w$ in period $p$, i.e.
\[
Y_{w,p,(i,j)} = 0\implies H_{w,p,(i,j)}\neq1.
\]
This constraint can be translated in a linear inequality by imposing
\[
H_{w,p,(i,j)} \leq Y_{w,p,(i,j)}.
\]

\paragraph{Main constraints.}


\textbf{(C1) Each pair of teams plays exactly once:}

\[
\forall (i,j) \in L \qquad
\sum_{w=1}^{W} \sum_{p=1}^{P} Y_{w,p,(i,j)} = 1.
\]

\textbf{(C2) Each team plays exactly once per week:}

\[
\forall t,w \qquad
\sum_{p=1}^{P}
\sum_{\substack{(i,j)\in L \\ i=t \vee j=t}}
Y_{w,p,(i,j)} = 1.
\]

\textbf{(C3) Each team plays at most twice in the same period:}

\[
\forall t,p \qquad
\sum_{w=1}^{W}
\sum_{\substack{(i,j)\in L \\ i=t \vee j=t}}
Y_{w,p,(i,j)} \leq 2.
\]

\textbf{(C4) Each slot $(w,p)$ hosts exactly one game:}

\[
\forall w,p \qquad
\sum_{(i,j)\in L} Y_{w,p,(i,j)} = 1.
\]

With this model, we do not need to explicitly impose that exactly two teams
play in each slot, since the possible games have already been represented as
team pairs in $L$.

\paragraph{Symmetry breaking constraints}
\begin{itemize}
    
    \item[(C5)] We fix each match in the first week by imposing:
    \begin{align*}
        \forall p =1,\dots,P\qquad   &Y_{1,p,(2p-1,\,2p)}=1\\
        &H_{1,p,(2p-1,\,2p)}=1
    \end{align*}
    
    \item[(C6)] Since the weeks are interchangeable in our formulation (constraints and objective do not depend on the week index), we can break the week-permutation symmetry by fixing the opponent of team 1 in each week, without fixing the period.
    \[
    \forall w = 1,\dots,W \qquad \sum_{p = 1}^P Y_{w,p,(1,\,w+1)} = 1,
    \]
    i.e., in week \(w\) team \(1\) plays against team \(w+1\) in some period \(p\).

    \item[(C7)] In addition we break symmetry by fixing the main diagonal of the schedule matrix: in each diagonal slot \((w,p)=(p,p)\) (for the first \(P\) weeks), team \(1\) must be scheduled, without fixing the opponent. Formally:
    \[
    \forall p = 1,\dots,P \qquad
    \sum_{j=2}^{T} Y_{p,p,(1,j)} = 1.
    \]
    
\end{itemize}

% ======= VALIDATION MIP =======
\subsection{Validation}

\paragraph{Experimental setup.}
All experiments for the MIP formulation were conducted using the CBC solver
through the \texttt{python-mip} library~\cite{pythonmip}. The implementation is
provided in Python, and the solver is invoked through the \texttt{python-mip}
API.

For each instance, two versions of the model were executed:
\begin{itemize}
    \item \texttt{mip\_cbc\_decision}, corresponding to the decision version of
    the problem, where the goal is to find any feasible schedule;
    \item \texttt{mip\_cbc\_optimization}, corresponding to the optimization
    version, where the objective is to minimize the total home/away imbalance.
\end{itemize}

The runtime is measured as total wall-clock time, including both model
construction and solving time. A global time limit of 300 seconds is imposed.
For solved instances, the reported time is the floor of the actual total runtime.
If the decision version does not find a feasible schedule within the time limit,
the runtime is reported as N/A. Similarly, for the optimization version, the
runtime is reported only when optimality is proved within the time limit.

For optimization runs that do not prove optimality within the time limit, the
objective value, when available, reports the best solution quality found by the
solver before timeout.

\paragraph{Experimental results.}

\begin{table}[H]
\centering
\begin{tabular}{c|cc|ccc}
\hline
& \multicolumn{2}{c|}{Decision}
& \multicolumn{3}{c}{Optimization} \\
\cline{2-6}
$n$ & Solved & Time (s)
& Optimal & Time (s) & Obj. \\
\hline
2  & Yes & 0   & Yes & 0   & 2  \\
4  & UNSAT & N/A & UNSAT & N/A & N/A \\
6  & Yes & 0   & Yes & 0   & 6  \\
8  & Yes & 0   & Yes & 0   & 8  \\
10 & Yes & 6   & Yes & 13  & 10 \\
12 & Yes & 87  & Yes & 270 & 12 \\
\hline
\end{tabular}
\caption{Experimental results of the MIP model using CBC. For decision runs,
``Solved'' indicates whether a feasible schedule was found within 300 seconds.
For optimization runs, ``Optimal'' indicates whether optimality was proved within
300 seconds. The runtime is reported as N/A when no feasible or optimal solution
is available for the corresponding setting. Objective values refer to the
home/away imbalance minimized by the optimization model.}
\label{tab:mip_results}
\end{table}
